\newpage
\section{Análise do Texto}

Nesta seção, serão apresentados comentários e reflexões acerca de \textit{Os
Lusíadas}, considerando tanto seus aspectos literários quanto os sentidos
presentes em sua construção poética.

\subsection{Considerações Preliminares}

Antes de analisarmos diretamente o uso da tradição clássica por Camões, convém
estabelecer algumas considerações sobre a literatura pagã, o mito e a religião
dos povos antigos. A partir delas, será possível compreender com maior clareza
como a tradição clássica é recebida e empregada na obra portuguesa.

\subsubsection{A Literatura Pagã à Luz de São Basílio}

\textit{Os Lusíadas} é epopeia por excelência, fruto do Renascimento. Composta
em dez cantos, organiza-se em oitavas de versos decassílabos e retoma a
tradição épica clássica, tanto em sua forma quanto em seus recursos poéticos.
Por isso, utiliza-se da mitologia greco-romana por toda sua extensão. E, diante
disso, surge uma questão de consciência para nós, enquanto católicos: ``Isto
não vai de encontro à minha fé?''

São Basílio Magno, em sua \citetitle{basilio_carta_jovens}
\cite{basilio_carta_jovens}, oferece-nos um princípio útil para responder tal
questão:

\begin{quote}
  ``Devemos estudar os autores profanos que louvam a virtude e condenam o
  vício. [...] Do mesmo modo, aqueles que nas suas leituras não procuram
  unicamente o prazer poderão extrair coisas muito úteis para o espírito.
  Enfim, devemos imitar as abelhas, que não voam indiscriminadamente sobre
  todos os tipos de flores: escolhem somente as mais aptas para lhes fornecer
  aquilo que é útil ao seu trabalho, abandonando o resto inútil. Faremos o
  mesmo se formos sábios: após colher nesses livros tudo aquilo que é precioso
  para o conhecimento da verdade, abandonaremos o resto.''
\end{quote}

Assim, a presença da mitologia pagã em \textit{Os Lusíadas} não precisa ser,
por si mesma, motivo para rejeitarmos a obra. Seguindo a recomendação de São
Basílio podemos aproveitar o que há de virtuoso na obra. Sem assim ferir a
moral cristã que prometemos seguir no Batismo.

\subsubsection{O Mito e a Transmissão do Conhecimento}

Para além dessa recomendação, à qual o leitor pode ou não aderir, há ainda
outras reflexões que podemos fazer. Seria um erro olhar para os povos antigos
e enxergar em seus mitos apenas histórias sobre deuses e heróis, como se estes
fossem tão somente personagens dotados de poderes sobrenaturais. O mito, embora
nascido no seio da religião e da tradição, pode carregar consigo algo que
ultrapassa a própria narrativa: uma maneira de contemplar e compreender a
ordem do mundo. Aquilo que não se escreve em fórmulas pode ser guardado em
imagens; aquilo que não se explica em conceitos pode ser transmitido em
símbolos.

É justamente nessa possibilidade que Giorgio de Santillana e Hertha von Dechend
se debruçam em \citetitle{santillana_dechend_hamlets_mill}
\cite{santillana_dechend_hamlets_mill}. Os autores procuram nos antigos mitos
os vestígios de uma linguagem celeste, na qual deuses, heróis e acontecimentos
poderiam conservar conhecimentos acerca da ordem e dos grandes ciclos dos céus.
O mito seria, assim, não apenas narrativa, mas também memória: uma forma
simbólica de transmitir conhecimentos astronômicos através das gerações,
especialmente aqueles relacionados à precessão dos equinócios.

\subsubsection{A Religião Antiga}

Além disso, a própria religião tradicional dos gregos não permaneceu
inteiramente imóvel diante das transformações intelectuais do período clássico.
O desenvolvimento da filosofia, sobretudo em autores como Platão e Aristóteles,
introduziu novas formas de pensar o divino, o mundo e sua ordem, tornando menos
simples a identificação entre a religião vivida e as narrativas tradicionais
sobre os deuses.

Uma distinção semelhante entre diferentes dimensões da religião antiga aparece
em Fustel de Coulanges. Em \citetitle{fustel_cidade_antiga}
\cite{fustel_cidade_antiga}, o autor distingue a religião pública da religião
doméstica. A primeira estava ligada aos deuses da cidade, aos ritos públicos e
às instituições cívicas; a segunda, ao culto dos antepassados e à vida da
família. Não se tratava, portanto, de duas expressões idênticas de uma mesma
religiosidade, mas de esferas distintas, ambas profundamente integradas à vida
antiga.

Uma analogia moderna pode tornar essa distinção mais clara. Alguém pode
celebrar o Natal, conservar costumes de origem cristã e reconhecer-se
culturalmente ligado ao cristianismo sem, contudo, frequentar a Missa ou
participar regularmente de sua vida religiosa. A presença de elementos
religiosos na cultura, portanto, não implica necessariamente uma adesão
integral à prática religiosa.

Algo semelhante deve ser considerado ao tratarmos da mitologia greco-romana.
A presença dos deuses nas narrativas, na poesia e nos costumes não permite
concluir, por si só, que cada grego ou romano concebesse literalmente tais
figuras como seres sobrenaturais da maneira como uma leitura moderna poderia
sugerir. Entre o mito narrado, o culto praticado e a compreensão intelectual
do divino havia uma realidade muito mais complexa. É justamente nesse espaço
que a tradição mítica pode ser compreendida não apenas como religião, mas
também como linguagem simbólica, literária e cultural.

\subsubsection{Camões e a Epopeia dos Portugueses}

É à luz dessas considerações que podemos compreender o uso da tradição clássica
por Camões. O poeta serve-se da literatura pagã e da erudição clássica como
ponte para cantar o herói coletivo: a própria pátria portuguesa. Seu intuito é
elevar os feitos históricos dos portugueses à dignidade da epopeia, contrapondo
a verdade histórica de suas conquistas às fábulas dos antigos. Assim, aquilo
que os poetas da Antiguidade haviam cantado por meio de mitos e deuses, Camões
procura cantar por meio dos feitos reais dos portugueses.

\newpage

\subsection{Sentido e Interpretação}

Já nas primeiras estrofes, Camões apresenta o canto como instrumento de
preservação da memória. Ao cantar as ``memórias gloriosas'' dos reis e daqueles
que, ``por obras valerosas'', se vão ``da lei da Morte libertando'', o poeta
afirma a capacidade da poesia de perpetuar os grandes feitos para além da vida.

A poesia assume, assim, uma forma de imortalidade terrena: não pela conservação
do corpo, mas pela permanência da memória. Ao cantar os feitos dignos de
lembrança, o poeta os retira do esquecimento e os transmite às gerações
seguintes, tornando-se guardião da memória lusitana.

É lamentável, contudo, que uma obra concebida para conservar a memória de um
povo já não ocupe, em nossa formação escolar, o espaço que sua importância
poderia sugerir. Embora \textit{Os Lusíadas} ainda encontre lugar no ensino
brasileiro, sua presença muitas vezes ocorre por meio de excertos, tornando
menos comum o contato integral do estudante com a obra. Cabe, assim, muitas
vezes, ao próprio leitor buscá-la e descobrir, por iniciativa própria, a
riqueza de uma tradição que deveria ser transmitida entre as gerações. A obra
corre o risco de ser reduzida a uma etapa necessária à preparação para o
vestibular, em vez de ser lida e apreciada como aquilo que é: uma grande obra
da literatura em língua portuguesa.

\ornamentalbreak

Na terceira estrofe, Camões estabelece uma comparação com a tradição épica da
Antiguidade e introduz o repertório mitológico que acompanhará \textit{Os
Lusíadas} até o fim. As navegações de Gregos e Troianos e as vitórias de
Alexandre e Trajano são tomadas como medida para revelar a grandeza dos feitos
portugueses. Ao ordenar que ``cesse tudo o que a Musa antiga canta'', o poeta
não rejeita os antigos, mas utiliza sua grandeza já consagrada para afirmar a
superioridade do novo canto.

É também nessa estrofe que aparecem Netuno e Marte, figuras mitológicas que
passarão a integrar a linguagem poética da epopeia. A expressão ``a quem Netuno
e Marte obedeceram'' associa essas divindades ao domínio do mar e da guerra,
exaltando os portugueses justamente nos campos em que seus feitos se
destacaram. A ``Musa antiga'' é, então, incorporada e superada: Camões emprega
a mitologia clássica como recurso poético para cantar os feitos históricos dos
portugueses.

\ornamentalbreak

Nas estrofes 4 e 5, Camões invoca as Tágides, ninfas do rio Tejo, pedindo-lhes
inspiração para elevar seu canto à altura dos feitos que pretende narrar. O
pedido manifesta a humildade do poeta diante da grandeza de sua matéria: para
cantar feitos tão sublimes, necessita de um canto igualmente grandioso. Pede,
assim, que sua poesia espalhe pelo mundo a memória da ``famosa gente'' e dê aos
feitos portugueses a dimensão universal que lhes atribui.

\ornamentalbreak

Nas estrofes 6 a 18, Camões dedica o poema a D. Sebastião, então jovem rei de
Portugal (nomeado em 1557, aos 3 anos), apresentando-o como esperança da
Cristandade e continuador da grandeza portuguesa. A dedicatória assume também
um caráter de exortação: o poeta conclama o monarca a prosseguir as conquistas
portuguesas e a combater os inimigos da fé, sobretudo os muçulmanos.

Camões recorre à linguagem épica e mitológica para exaltar o rei. Tétis,
``afeiçoada ao gesto belo e tenro'' de D. Sebastião, é apresentada como
desejosa de tê-lo por genro, dispondo para ele todo o seu domínio marítimo como
dote. O poeta afirma ainda que seus inimigos, ao vê-lo, já demonstram temor e
se inclinam ao seu jugo, reforçando a imagem de um monarca destinado à grandeza
e à expansão do poder português.

O poeta estabelece ainda uma comparação decisiva entre os heróis da Antiguidade
e os portugueses. \textit{Nuno Álvares Pereira}, condestável e célebre
comandante militar português, \textit{Egas Moniz}, aio de D. Afonso Henriques,
e \textit{D. Fuas Roupinho}, nobre e guerreiro do século XII, são apresentados
como exemplos da grandeza histórica de Portugal. Aos \textit{Doze Pares de
França}, lendários companheiros de Carlos Magno, Camões contrapõe os
\textit{Doze de Inglaterra}, grupo de cavaleiros portugueses que, segundo a
tradição, foram à Inglaterra defender a honra de doze damas, destacando entre
eles \textit{Magriço}. Por fim, \textit{Vasco da Gama}, comandante da expedição
que alcançou a Índia por via marítima, é comparado a \textit{Eneias}, herói
troiano celebrado na epopeia de Virgílio.

Assim, Camões coloca lado a lado personagens históricas, lendárias e míticas,
mas estabelece entre elas uma diferença fundamental: enquanto a Antiguidade
celebrava feitos revestidos pela fábula, os portugueses possuem feitos
históricos cuja grandeza dispensa a invenção.

A história, contudo, daria à dedicatória um significado trágico. Em 4 de agosto
de 1578, D. Sebastião desapareceu na Batalha de Alcácer-Quibir, sem que seu
corpo fosse identificado. O episódio alimentaria posteriormente o
\textit{sebastianismo}, crença na volta do rei desaparecido para restaurar
Portugal em seus momentos de maior necessidade
\cite{profetismo_espiritualidade_camoes_pascoaes}. Nesse imaginário, D.
Sebastião aproximou-se de outras figuras europeias do rei desaparecido ou
adormecido, como o Rei Artur, Frederico Barbarossa e Constantino XI Paleólogo
\cite{zierer_rei_artur_sebastiao}.

A própria trajetória posterior de Camões acrescenta uma dimensão simbólica à
relação entre o poeta e a pátria que cantou. Após a derrota de Alcácer-Quibir e
a incorporação de Portugal à Monarquia Hispânica, a morte de Camões, em 1580,
parece encerrar simbolicamente o mesmo ciclo de crise que atingia a pátria que
cantara: o poeta desaparece justamente quando Portugal perde sua autonomia. A
localização de seus restos mortais permanece incerta, e seu possível
desaparecimento após o terremoto de Lisboa de 1755 torna ainda mais nebuloso o
destino de seus restos \cite{garrett_camoes}. Poeta e pátria parecem, assim,
compartilhar um mesmo destino: aquele que se fizera voz da memória portuguesa
desaparece no momento em que a própria pátria atravessa uma de suas maiores
crises. A imagem aproxima-se da própria concepção camoniana da poesia: se o
corpo desaparece, resta o canto para preservar aquilo que a morte e o tempo não
deveriam apagar.
